% Chapter: 06_discussion
% Included from lectures/main.tex

{
\setbeamertemplate{section page}{
  \begin{centering}
    \usebeamerfont{section title}\insertsection\par
  \end{centering}
  \note{
  \begin{itemize}
  \item You now have a rough overview of what happened the last years in MechInterp
  \item Where do we go from here?
  \item Open question: what kind of progress would you want to see in MechInterp? What is realistic?
  \end{itemize}
  }
}
\section{Discussion}
}

% \problemcard{x}{y}{title}{bullets}
\newcommand{\problemcard}[4]{\node[card,anchor=north west,text width=6.4cm,minimum height=2.5cm] at (#1,#2) {\textbf{#3}\\[2pt]{\footnotesize#4}};}

\begin{frame}{Open Problems: Methods}
\centering
\begin{tikzpicture}[x=1cm,y=1cm]
\problemcard{0}{3}{Decomposition}{$\bullet$ SAEs have high reconstruction errors\\$\bullet$ Expensive for large models\\$\bullet$ Sparsity is a flawed proxy for interpretability\\$\bullet$ No agreed definition of a ``feature''}
\problemcard{7.3}{3}{Description}{$\bullet$ Top activating examples are correlational\\$\bullet$ Attribution is a first-order approximation\\$\bullet$ Causal interventions are expensive at scale}
\problemcard{3.65}{-0.1}{Validation}{$\bullet$ Different methods give different explanations\\$\bullet$ Rarely tested on real-world tasks\\$\bullet$ Hypotheses get conflated with conclusions}
\end{tikzpicture}

\source{sharkey2025}

\note{
\begin{itemize}
\item Neurons, attention heads, and layers are polysemantic and don't decompose cleanly
\item SDL (SAEs etc.) has high reconstruction errors, is expensive for large models, assumes linearity in nonlinear models, and uses sparsity as a flawed proxy for interpretability
\item SDL latents are dataset-dependent and may not contain the concepts needed for downstream use
\item Highly activating examples are correlational and prone to interpretability illusions
\item Attribution methods are first-order approximations and can be adversarially fooled
\item Different methods yield different interpretations for the same phenomenon; lack of standard benchmarks and model organisms
\end{itemize}
}
\end{frame}

\begin{frame}{Open Problems: Scaling and Applications}
\centering
\begin{tikzpicture}[x=1cm,y=1cm]
\problemcard{0}{3}{Automation}{$\bullet$ Fully automating current pipelines would not give satisfying explanations\\$\bullet$ ``Streetlight interpretability'': biased towards simple tasks}
\problemcard{7.3}{3}{Applications}{$\bullet$ Can't reliably separate features that \emph{recognize} deception from mechanisms that \emph{cause} it\\$\bullet$ Mechanistic verification far beyond current capabilities\\$\bullet$ Steering and unlearning limited by imprecise mechanisms}
\end{tikzpicture}

\source{sharkey2025}

\note{
\begin{itemize}
\item Circuit discovery has low faithfulness, struggles with backup/negative behavior, and is biased toward simple tasks
\item Monitoring/auditing: can't yet reliably distinguish features that recognize deception from mechanisms that cause deceptive behavior
\item Control: activation steering, unlearning, and model editing are all limited by the inability to isolate individual mechanisms
\item Predicting behavior: formal verification of AI is well beyond current capabilities
\item Positive side: mechanistic understanding could also enable pruning, distillation, data selection, and microscope AI for science
\end{itemize}
}
\end{frame}

\begin{frame}{Discussion}
\begin{itemize}
\setlength{\itemsep}{0.8em}
\item Which critique from the first hour do you find most convincing? Why?
\item What result would convince you that interpretability can rule out deception?
\item Ambitious reverse-engineering or pragmatic tools: where would you work?
\end{itemize}

\vspace{1.5em}
{\small\textbf{Further reading:}}
\begin{enumerate}
\small
\item \emph{A Pragmatic Vision for Interpretability} (Nanda et al., 2025)
\item \emph{Against Almost Every Theory of Impact of Interpretability} (Ségerie, 2023)
\item \emph{The Misguided Quest for Mechanistic AI Interpretability} (Hendrycks, 2025)
\item \emph{Activation Space Interpretability May Be Doomed} (Chughtai, 2025)
\end{enumerate}

\note{
\begin{itemize}
\item Open discussion to close the session
\end{itemize}
}
\end{frame}
